\documentclass[11pt,a4paper]{article}
\usepackage[T1]{fontenc}
\usepackage[utf8]{inputenc}
\usepackage{lmodern}
\usepackage[margin=27mm,headheight=15pt]{geometry}
\usepackage{amsmath,amssymb,amsthm,mathtools}
\usepackage{booktabs,longtable,array,enumitem,microtype}
\usepackage{xcolor,fancyhdr}
\usepackage[unicode,colorlinks=true,linkcolor=blue!45!black,citecolor=blue!45!black,urlcolor=blue!45!black]{hyperref}
\usepackage{bookmark}
\definecolor{ink}{RGB}{25,42,64}
\definecolor{accent}{RGB}{35,91,120}
\newcommand{\R}{\mathbb R}
\newcommand{\N}{\mathbb N}
\newcommand{\E}{\mathbb E}
\newcommand{\Pbb}{\mathbb P}
\newcommand{\norm}[1]{\left\lVert#1\right\rVert}
\newcommand{\abs}[1]{\left\lvert#1\right\rvert}
\newcommand{\dd}{\,\mathrm d}
\DeclareMathOperator{\dist}{dist}
\DeclareMathOperator{\supp}{supp}
\DeclareMathOperator{\spanop}{span}
\DeclareMathOperator{\diam}{diam}
\DeclareMathOperator{\Lip}{Lip}
\newtheorem{theorem}{Theorem}[section]
\newtheorem{lemma}[theorem]{Lemma}
\newtheorem{proposition}[theorem]{Proposition}
\newtheorem{corollary}[theorem]{Corollary}
\theoremstyle{definition}
\newtheorem{definition}[theorem]{Definition}
\theoremstyle{remark}
\newtheorem{remark}[theorem]{Remark}
\numberwithin{equation}{section}
\setlist{topsep=4pt,itemsep=2pt,parsep=0pt}
\setlength{\emergencystretch}{2em}
\pagestyle{fancy}
\fancyhf{}
\fancyhead[L]{\small\color{ink}AIM 523}
\fancyhead[R]{\small 4 November 2026}
\fancyfoot[C]{\thepage}
\renewcommand{\headrulewidth}{0.3pt}


\urlstyle{same}
\title{An exact error floor for adaptive Monte Carlo approximation of a noncompact operator}
\author{Matthew J. Colbrook\thanks{The solutions were found by AI. The submitted arguments and the repository audit are AI work; independent human review and formal verification are not claimed.}\\[4pt]
\small Department of Applied Mathematics and Theoretical Physics\\
\small University of Cambridge, Cambridge CB3 0WA, United Kingdom\\
\small\href{mailto:m.colbrook@damtp.cam.ac.uk}{m.colbrook@damtp.cam.ac.uk}}
\date{4 November 2026}
\hypersetup{pdftitle={An exact error floor for adaptive Monte Carlo approximation of a noncompact operator},pdfauthor={Matthew J. Colbrook},pdfsubject={AIM 523: AI-found solution or partial result}}
\begin{document}
\maketitle
\begin{abstract}
For the canonical embedding $\ell_1\hookrightarrow\ell_\infty$, we prove that the optimal worst-case root-mean-square error is exactly $1/2$ at every positive finite expected information budget. The proof permits adaptive continuous linear information and the joint measurability conventions of AIM 523. We also give compactness and separable-target consequences.
\end{abstract}
\noindent\textbf{Catalogue target.} \href{https://github.com/MColbrook/AIM/blob/c4036871b0bf15cbc83e79c780b7191bcbb2cfed/problems/523-can-adaptive-monte-carlo-approximate-a-noncompact-operator.md}{AIM 523 at repository revision \texttt{c4036871b0bf}}. The package preserves its full statement in \texttt{statement.md}; the review and target comparison are linked from the package README.
\medskip

\section{Problem 523: an exact error floor for adaptive Monte Carlo}
\label{sec:523}

Let \(I:\ell_1\to\ell_\infty\) be the identity inclusion. The input set is
\(B=\{x:\norm{x}_1<1\}\). Measurements are continuous linear functionals,
their choice and the stopping rule may depend arbitrarily on the preceding
answers and the random seed, and the output and cardinality maps have the
joint measurability required in the problem. The cost bound is
\(\sup_{x\in B}\E N(x,\omega)\le n\), and the error is the worst-input
root mean square error. More explicitly, the seed ranges over an arbitrary probability space, $N(x,\omega)$ is finite and takes values in $\N\cup\{0\}$, and $(x,\omega)\mapsto N(x,\omega)$ and $(x,\omega)\mapsto A_\omega(x)$ are jointly measurable for the norm Borel structures and the seed sigma algebra. We write
\[
e_n(I)=\inf_A\sup_{x\in B}\bigl(\E_\omega\|x-A_\omega(x)\|_\infty^2\bigr)^{1/2},
\]
where the infimum is over those algorithms satisfying the stated expected cost bound.

\begin{theorem}
\label{thm:523}
In this model,
\[
                 \boxed{e_n(I)=\frac12\qquad(n\ge1).}
\]
\end{theorem}

The issue left open in the sources is adaptivity, rather than the
nonadaptive error floor; see \cite{hnw2011,knw2024}. The proof below treats arbitrary
adaptive decision rules. Its main ingredient is a tightness statement for
the actual measurable output, which avoids imposing measurability on the
rules selecting the queries.

\subsection{A tightness lemma for finite adaptive information}

\begin{lemma}[Uniform tightness under bounded shifts]
\label{lem:tightness}
Let \(E\) be a separable real Banach space, and let \(\mu\) be a centered
Radon Gaussian measure on \(E\), with full support and finite second
moment. Let \(U\subset E\) be open. Suppose that
\(R:U\to\N_0\) is Borel measurable and is the output of an algorithm using
at most \(k\) adaptive continuous affine scalar measurements.

Let \(V\subset E\) be countable and bounded, and let
\(w:E\to[0,\infty)\) be bounded and Borel measurable. Assume that, for
some \(\eta>0\),
\[
 \operatorname{dist}(x+v,E\setminus U)\ge\eta
 \quad\text{whenever }w(x)>0,\ v\in V.
 \tag{T1}
\]
Extend \(R\) by zero outside \(U\). Then
\[
 \lim_{M\to\infty}\sup_{v\in V}
       \int_E w(x)\mathbf1_{\{R(x+v)\ge M\}}\,d\mu(x)=0.
 \tag{T2}
\]
Neither the query-selection rules nor the decoder need be measurable as
functions of the answer history.
\end{lemma}

\begin{proof}
We induct on \(k\). With no measurements, the output is a fixed integer
on \(U\); (T1) then proves the assertion. An algorithm that stops early
can be padded with zero measurements.

The first affine query has a fixed continuous linear part \(a\in E'\).
Its known additive constant may be subtracted from the answer. If
\(a=0\), freeze this first answer and apply induction to the remaining
queries. Otherwise, for \(Z\sim\mu\), put
\[
 \sigma^2=\E[a(Z)^2]>0,\qquad
 c=\E[Z\,a(Z)]\in E,\qquad
 m(t)=\frac{c}{\sigma^2}t,\qquad E_0=\ker a .
 \tag{T3}
\]
The Bochner expectation exists by the second-moment assumption.
Moreover \(a(c)=\sigma^2\), so \(a(m(t))=t\). Gaussian regression gives
\[
              Z=m(Y)+Z_0,\qquad Y=a(Z),
 \tag{T4}
\]
where \(Y\) has the \(N(0,\sigma^2)\) density \(p\), and \(Z_0\) is an
independent centered Gaussian random element of \(E_0\).
Indeed, every continuous linear functional of
\(Z-m(a(Z))\) has zero covariance with \(a(Z)\); joint Gaussianity
then gives independence. The law \(\mu_0\) of \(Z_0\) has full support
in \(E_0\), because
\[
                   P=\operatorname{Id}-m\circ a:E\to E_0
\]
is a continuous surjection and \(\mu\) has full support.
Its second moment is finite. Thus the Gaussian hypotheses are preserved
under conditioning.

For \(v\in V\), write
\[
                  b=a(v),\qquad v'=v-m(b)\in E_0.
 \tag{T5}
\]
The scalar \(b\) and the residual shift \(v'\) range over bounded sets.
Define the nonnegative integrable Borel function
\[
                  f(y,z)=p(y)w(m(y)+z)
       \quad\text{on }\R\times E_0.
\]
Changing the first-answer variable from \(t\) to \(y=t+b\) gives
\[
\begin{split}
 I_M(v)&:=\int_Ew(x)\mathbf1_{\{R(x+v)\ge M\}}\,d\mu(x)\\
 &=\int_{\R}\int_{E_0}
       f(y-b,z)\,
       \mathbf1_{\{R(m(y)+z+v')\ge M\}}\,d\mu_0(z)\,dy.
\end{split}
 \tag{T6}
\]

Fix \(\varepsilon>0\). Translation in the first real variable is
continuous in \(L^1(dy\,d\mu_0)\). Choose \(h>0\) such that
\[
 \norm{f(\,\cdot-t,\cdot)-f}_{L^1}<\varepsilon
 \quad\text{and}\quad \norm{m(t)}_E<\eta/2
                   \qquad(\abs t<h).
 \tag{T7}
\]
Since \(a(V)\) is bounded, finitely many intervals of radius less than
\(h\), with centers \(b_1,\ldots,b_L\), cover it. Partition \(V\) into
corresponding countable sets \(V_l\), and put
\[
                 W_l=\{v-m(a(v)):v\in V_l\}\subset E_0 .
 \tag{T8}
\]
Each \(W_l\) is countable and bounded. For \(v\in V_l\), replacing
\(f(y-a(v),z)\) in (T6) by \(f(y-b_l,z)\) changes the integral by at
most \(\varepsilon\), uniformly in \(v\) and \(M\).

Fix \(l\) and \(y\), and define
\[
\begin{split}
 U_y&=(U-m(y))\cap E_0,\\
 R_y(z)&=R(m(y)+z)\quad(z\in U_y),\\
 w_{l,y}(z)&=w(m(y-b_l)+z).
\end{split}
 \tag{T9}
\]
The function \(R_y\) is Borel measurable. All inputs \(m(y)+z\),
\(z\in E_0\), have the same first answer. Consequently \(R_y\) uses at
most \(k-1\) adaptive continuous affine measurements, with that answer
frozen.

The shifted-support condition remains valid. If \(w_{l,y}(z)>0\), let
\(x=m(y-b_l)+z\). For \(v\in V_l\), with \(v'\) as in (T5),
\[
             m(y)+z+v'=x+v+m(b_l-a(v)).
 \tag{T10}
\]
By (T1) and (T7), this point has distance at least \(\eta/2\) from
\(E\setminus U\). Distance within the affine subspace
\(m(y)+E_0\) is at least ambient distance. Hence
\[
 \operatorname{dist}_{E_0}(z+v',E_0\setminus U_y)\ge\eta/2
 \quad(w_{l,y}(z)>0,\ v'\in W_l).
 \tag{T11}
\]
Notice that the unshifted point \(m(y)+z\) need not lie in \(U\).
This is why the hypothesis was formulated as (T1), rather than as a
condition on the unshifted support alone.

Induction now implies, for every fixed \(y\),
\[
 T_{M,l}(y):=
 \sup_{v'\in W_l}\int_{E_0}w_{l,y}(z)
       \mathbf1_{\{R_y(z+v')\ge M\}}\,d\mu_0(z)
                  \longrightarrow0.
 \tag{T12}
\]
The supremum is countable. Each term is Borel measurable in \(y\),
directly from the Borel measurability of the original output \(R\)
and of \(w\). Thus \(T_{M,l}\) is Borel, and
\(0\le T_{M,l}\le\norm w_\infty\). Dominated convergence gives
\[
                 \int_\R p(y-b_l)T_{M,l}(y)\,dy\longrightarrow0 .
 \tag{T13}
\]
There are only finitely many \(l\). Combining (T6)--(T13) yields
\[
                    \limsup_{M\to\infty}\sup_{v\in V}I_M(v)
                         \le\varepsilon .
\]
Letting \(\varepsilon\downarrow0\) proves (T2). No history-to-query
map has been integrated.
\end{proof}

\subsection{Proof of the exact error floor}

\begin{proof}[Proof of Theorem~\ref{thm:523}]
For the upper bound, one measurement suffices:
\[
 s=\sum_{j\ge1}x_j,\qquad A(x)=(s/2,s/2,\ldots).
\]
This is an admissible \(\ell_\infty\)-valued output, and
\[
 \abs{x_j-s/2}
   =\frac12\abs{x_j-\sum_{i\ne j}x_i}
   \le\frac12\norm x_1 .
 \tag{523.1}
\]

For the lower bound, first impose a strict finite budget \(k\).
Let
\[
                    Z_j=2^{-j}g_j,
\]
where the \(g_j\) are independent standard Gaussians.
Then \(Z\in\ell_1\) almost surely and \(\E\norm Z_1^2<\infty\).
Its law has full support in \(\ell_1\): a finite Gaussian head can
approximate any prescribed finite vector with positive probability,
and the independent tail can be made small with positive probability
by first making its expected norm small.

Choose
\[
        0<r<1,\qquad 0<\rho<\min\{r/2,1-r\},\qquad
        \tau=r/2-\rho>0 .
 \tag{523.2}
\]
Condition \(Z\) on \(\norm Z_1\le\rho\), an event of positive
probability. Equivalently, use the bounded weight
\[
 w(z)=\frac{\mathbf1_{\{\norm z_1\le\rho\}}}
                 {\Pbb(\norm Z_1\le\rho)}.
 \tag{523.3}
\]
Take \(U=B\) and \(V=\{re_j:j\ge1\}\). The margin in
Lemma~\ref{lem:tightness} is at least \(1-r-\rho>0\).

For a fixed seed \(\omega\), define from the algorithm's output
\[
 R_\omega(x)=\min\{i\ge1:(A_\omega(x))_i>r/2\},
 \tag{523.4}
\]
with value zero if the set is empty. Coordinate evaluation is continuous
on \(\ell_\infty\), so \(R_\omega\) is Borel. It uses the same
measurements as \(A_\omega\). If \(x=z+re_j\), \(\norm z_1\le\rho\),
and \(\norm{x-A_\omega(x)}_\infty<\tau\), then
\[
 (A_\omega(x))_j>r/2,\qquad
 (A_\omega(x))_i<r/2\quad(i\ne j).
 \tag{523.5}
\]
Thus \(R_\omega(x)=j\). Uniform tightness implies
\[
 \Pbb\bigl(R_\omega(Z+re_j)=j\mid\norm Z_1\le\rho\bigr)
       \longrightarrow0,
 \tag{523.6}
\]
because this probability is bounded by the uniformly vanishing tail
at threshold \(j\).

The probabilities in (523.6) are measurable functions of \(\omega\),
by the joint measurability of the actual output. They are bounded by
one. Dominated convergence therefore permits averaging (523.6) over
any seed probability space. The norm error is jointly measurable even
though \(\ell_\infty\) is nonseparable: the map \(x\mapsto Ix\) has
separable range, so approximate it by countably valued Borel maps and
take limits of their norm differences from the actual output.
If \(e\) is the uniform RMS error, every shifted input law in
(523.6) is supported in \(B\), and hence
\[
 e^2\ge
 \E_{Z,\omega}\bigl[\norm{Z+re_j-A_\omega(Z+re_j)}_\infty^2
                    \mid\norm Z_1\le\rho\bigr]
 \ge \tau^2(1-o(1)).
 \tag{523.7}
\]
Let \(j\to\infty\), then \(\rho\downarrow0\) and \(r\uparrow1\).
This proves \(e\ge1/2\) for every strict finite budget.

For an expected budget, truncate the algorithm after \(k\) queries
and return zero if it has not stopped. The truncated output is jointly
Borel, since the original output and cardinality are jointly Borel.
For every \(x\in B\),
\[
\begin{split}
 \E\norm{x-A^{(k)}_\omega(x)}_\infty^2
 &\le \E\norm{x-A_\omega(x)}_\infty^2
          +\norm x_\infty^2\,\Pbb(N(x,\omega)>k)\\
 &\le e^2+n/k .
\end{split}
 \tag{523.8}
\]
The strict-budget result gives \(e^2+n/k\ge1/4\) for every \(k\).
Letting \(k\to\infty\) completes the proof.
\end{proof}

\begin{corollary}[If the output is required to belong to \(c_0\)]
\label{cor:523-c0}
For the inclusion \(I:\ell_1\to c_0\), equipped with the same
continuous-linear information, joint output/cardinality measurability,
expected-budget convention, and worst-input RMS error,
\[
                    e_n(I:\ell_1\to c_0)=1
                 \qquad\text{for every finite }n.
\]
\end{corollary}
\begin{proof}
Fix \(\varepsilon>0\). For a deterministic realization with output
\(A(x)\in c_0\), the integer
\[
 R_A(x)=\max\{j\ge1:\abs{A(x)_j}>\varepsilon\},
\]
with value zero when the set is empty, is finite and Borel measurable.
Indeed, \(\{R_A\le M\}=\bigcap_{j>M}\{\abs{A(x)_j}\le\varepsilon\}\).
Use the same Gaussian law, conditioning, and shifts \(re_j\) as in
the preceding proof. If \(x=z+re_j\), \(\norm z_1\le\rho\), and
\[
                   \norm{x-A(x)}_\infty<r-\rho-\varepsilon,
\]
then \(\abs{A(x)_j}>\varepsilon\), and therefore \(R_A(x)\ge j\).
Lemma~\ref{lem:tightness} makes this event's conditioned
probability tend to zero, for each seed, as \(j\to\infty\).
Dominated convergence over the seed gives an RMS lower bound
\(r-\rho-\varepsilon\) for every strict finite query budget.
Let \(\rho,\varepsilon\downarrow0\) and \(r\uparrow1\).
For expected budgets, truncation and zero output on unfinished paths
give the same estimate \(e^2+n/k\) as in (523.8), so the lower
bound remains one. The zero algorithm attains one.
\end{proof}

\begin{remark}
The value \(1/2\) for the \(\ell_\infty\) codomain uses the constant
sequence output of the one-query algorithm. A nonzero constant
sequence does not belong to \(c_0\); restricting the allowed output
space therefore changes the exact randomized error.
\end{remark}

\subsection{A direct compactness criterion}

\begin{corollary}
\label{cor:compactness}
Let \(S:X\to Y\) be a bounded linear operator between real Banach
spaces, with the same continuous-linear information, joint
output/cardinality measurability, expected-budget convention, and
worst-input RMS error. Then
\[
                 \lim_{n\to\infty}e_n(S)=0
                    \quad\Longleftrightarrow\quad
                 S\text{ is compact}.
 \tag{523.9}
\]
More quantitatively, if there are \(v_j\in X\), \(\norm{v_j}_X\le1\),
such that \(\norm{Sv_i-Sv_j}_Y\ge\alpha>0\) for \(i\ne j\), then
\(e_n(S)\ge\alpha/2\) for every finite \(n\).
\end{corollary}

\begin{proof}
Suppose first that \(S\) is noncompact and choose such a separated
sequence. Its closed linear span \(E\subset X\) is a separable Banach
space. A full-support centered Gaussian with finite second moment
exists on \(E\). For example, if \((u_l)\) lies in its unit ball and
has dense linear span, the series
\[
                         Z=\sum_{l\ge1}2^{-l}g_lu_l
 \tag{523.10}
\]
converges absolutely almost surely and in \(L^2(E)\); the finite-head
and independent-tail argument used above proves full support.

Choose \(0<r<1\), and then choose \(\rho>0\) so that
\[
         r+\rho<1,\qquad \tau=r\alpha/2-\norm S\,\rho>0.
\]
Condition on \(\norm Z_X\le\rho\) and apply
Lemma~\ref{lem:tightness} with the shifts \(rv_j\). Classify the output
by
\[
 R_A(x)=\min\{j:\norm{A(x)-rSv_j}_Y<r\alpha/2\},
 \tag{523.11}
\]
with value zero when no such \(j\) exists. The balls in (523.11) are
disjoint, and this is a Borel integer output. If
\(x=Z+rv_j\) and \(\norm{Sx-A(x)}_Y<\tau\), then (523.11) returns
\(j\). The same tightness and averaging argument as in (523.6)--(523.7)
gives error at least \(\tau\). Let \(\rho\downarrow0\) and
\(r\uparrow1\).

Restriction to \(E\) preserves all required measurability. In
particular, \(S(E)\) is separable, so the norm error is jointly
measurable even when \(Y\) is nonseparable: approximate \(Sx\) by
countably valued Borel maps and take limits of the corresponding norm
differences. Expected budgets are handled without any projection of
the \(Y\)-valued output:
\[
 \E\norm{Sx-A^{(k)}(x)}_Y^2
       \le e^2+\norm S^2\,n/k .
 \tag{523.12}
\]
Thus the lower bound \(\alpha/2\) survives expected cardinality.

Conversely, suppose that \(K=\overline{S(B_X)}\) is compact. Fix
\(\varepsilon>0\). For each \(z\in K-K\) with \(\norm z_Y\ge\varepsilon\),
choose a norm-one functional \(y'_z\in Y'\) with
\(y'_z(z)=\norm z_Y\). Compactness gives finitely many of these
functionals, \(y'_1,\ldots,y'_m\), such that
\[
 z,w\in K,\quad
 \max_i\abs{y'_i(z-w)}\le\varepsilon/2
                 \quad\Longrightarrow\quad
 \norm{z-w}_Y<\varepsilon .
 \tag{523.13}
\]
Indeed, the open sets
\(\{q:\abs{y'_z(q)}>\varepsilon/2\}\) cover the compact subset
\(\{q\in K-K:\norm q_Y\ge\varepsilon\}\).

Measure the fixed functionals \(L_i=y'_i\circ S\), and take a finite
\(\varepsilon/4\)-net \(z_1,\ldots,z_N\) in \(K\). Given the data
\(t=(t_i)\), return the first \(z_j\) satisfying
\[
                  \max_i\abs{t_i-y'_i(z_j)}\le\varepsilon/4,
 \tag{523.14}
\]
and return zero outside the union of these finitely many boxes.
For actual data, such a center exists; (523.13) gives error less than
\(\varepsilon\). This decoder is Borel and has finite range.
If the compact set used in (523.13) is empty, a constant decoder
already suffices. Hence every desired accuracy is achieved by
deterministic nonadaptive continuous-linear information. No
approximation property of \(X\) or \(Y\) is needed.
\end{proof}

\begin{corollary}[A quantitative lower bound for separable codomains]
\label{cor:noncompactness-radius}
Under the hypotheses of Corollary~\ref{cor:compactness}, suppose also
that \(Y\) is separable. Define the Hausdorff noncompactness radius by
\[
 \chi(S(B_X))=\inf\bigl\{r>0:S(B_X)\text{ is covered by finitely
 many open }Y\text{-balls of radius }r\bigr\}.
\]
Then every finite expected budget satisfies
\[
                         e_n(S)\ge\chi(S(B_X)).
\]
\end{corollary}
\begin{proof}
Assume \(\chi(S(B_X))>0\), and fix \(0<r<\chi(S(B_X))\).
Choose a countable \(V\subset B_X\) such that \(S(V)\) is dense
in \(S(B_X)\), and let \(E\) be its closed linear span in \(X\).
Use a full-support Gaussian \(Z\) on \(E\), conditioned on
\(\norm Z_X\le\rho\). Fix \(0<\theta<1\) with \(\theta+\rho<1\),
and use the bounded countable shifts \(\theta V\).

Let \((y_i)\) be dense in \(Y\). For \(\varepsilon>0\), set
\[
       R_\varepsilon(y)=\min\{i\ge1:\norm{y-y_i}_Y<\varepsilon\}.
\]
This is a finite Borel integer. For a strict finite-query algorithm,
Lemma~\ref{lem:tightness} applied to \(R_\varepsilon\circ A_\omega\)
gives, for each seed,
\[
 q_M(\omega):=\sup_{v\in V}
 \Pbb_Z\!\left(R_\varepsilon(A_\omega(\theta v+Z))>M
                    \,\middle|\,\norm Z_X\le\rho\right)
                         \longrightarrow0.
\]
Countability of \(V\) and joint output measurability make \(q_M\)
measurable in the seed. Thus \(\E_\omega q_M\to0\).

For each \(M\), the definition of \(\chi\) and the density of \(S(V)\)
give \(v_M\in V\) with
\[
              \min_{i\le M}\norm{\theta Sv_M-y_i}_Y>\theta r.
\]
Indeed, failure would give a finite closed-ball cover of
\(S(B_X)\) with radius \(r\), and hence open-ball covers of every
larger radius, contrary to \(r<\chi\).
On \(R_\varepsilon(A_\omega(\theta v_M+Z))\le M\), the error is
at least \(\theta r-\norm S\rho-\varepsilon\). Its probability tends
to one by the uniform averaged bound above. The worst-input RMS
error is consequently at least this number. Let
\(\rho,\varepsilon\downarrow0\), then \(\theta\uparrow1\), and
finally \(r\uparrow\chi(S(B_X))\).

For expected budgets, truncate after \(k\) queries and use
(523.12). The strict-budget lower bound gives
\(e^2+\norm S^2n/k\ge\chi(S(B_X))^2\); let \(k\to\infty\).
\end{proof}


\begin{thebibliography}{99}
\bibitem{hnw2011}A. Hinrichs, E. Novak and H. Wo\'zniakowski, \emph{Discontinuous information in the worst case and randomized settings}, Math. Nachr. \textbf{286} (2013), 679--690. Primary manuscript: \url{https://arxiv.org/abs/1106.2945}; \url{https://doi.org/10.1002/mana.201100128}.
\bibitem{knw2024}R. J. Kunsch, E. Novak and M. Wnuk, \emph{Randomized approximation of summable sequences---adaptive and non-adaptive}, J. Approx. Theory \textbf{304} (2024), 106056. \url{https://arxiv.org/abs/2308.01705}; \url{https://doi.org/10.1016/j.jat.2024.106056}.
\end{thebibliography}
\end{document}
